\documentclass[12pt]{article}
\usepackage[T1]{fontenc}
\usepackage[margin=1in]{geometry}
\usepackage{amsmath,amssymb}
\usepackage{graphicx}
\usepackage{float}
\usepackage{array}
\usepackage{longtable}
\usepackage{booktabs}
\usepackage{threeparttable}
\usepackage{caption}
\usepackage{authblk}
\usepackage[authoryear,round]{natbib}
\usepackage{xurl}
\usepackage[hidelinks]{hyperref}
\usepackage{microtype}
\emergencystretch=1.5em
\renewcommand{\thesection}{\Roman{section}}
\renewcommand{\thesubsection}{\thesection.\Alph{subsection}}
% Possessive citation: one macro so a bibliography change cannot desynchronize
% hand-rolled \citeauthor{X}'s (\citeyear{X}) constructions.
\newcommand{\citepos}[1]{\citeauthor{#1}'s (\citeyear{#1})}

\title{Mortgage Lock-In and the Federal Reserve's Quantitative Tightening Shortfall}
\author{Eugene Ong}
\affil{Carnegie Mellon University}
\date{}
\begin{document}
\maketitle
\begin{abstract}
\noindent U.S. mortgages are fixed-rate, repaid at face value, and not portable, so rising rates lock households into old, cheap loans. Repayments on the Federal Reserve's mortgage bonds ran \$764.7 billion below its phased runoff ceilings over June 2022--November 2025, ceilings set roughly 1.7--1.9 times above the Federal Reserve's own contemporaneous projection of achievable runoff: the New York Fed's May 2022 staff baseline implied \$740.6 billion over the window, against the caps' \$1{,}417.5 billion \citep{nyfed2022}. That the caps would sit above achievable runoff is the staff's observation; I add its quantification. Most of the gap was never about the rate-responsive margin. A rate-inelastic baseline --- scheduled amortization plus a baseline turnover floor read from realized, partly behavioral turnover --- accounts for 85.7\% of the shortfall at the off-window floor and 88.7\% at the in-sample calibration point, both under the production floor form (35.6\% under the additive form). A loan-by-loan model with the response switched on accounts for 91.3\%, on the benchmark-consistent basis at the off-window floor this paper headlines. Measured against the projection instead, the unanticipated shortfall is \$87.8 billion; lock-in accounts for roughly half of it under the central allocation, and for 23--80\% across the disclosed floors and allocations, all upper bounds by construction. What lock-in itself adds is secondary, and a range rather than a number: $+2.3$ to $+9.1$ points under the production floor form --- the floor read's own sampling error at the central elasticity, wild-cluster on 31 clusters, the binding interval --- inside a form-conditional hull of $+3.5$ to $+13.1$, with $+5.6$ points, \$42.6 billion, an anchor convention rather than a central tendency. Measured corrections to the floor and the accounting basis move it in both directions, and its elasticity is imported, never re-measured here. The design identifies a marginal, not a level and not monthly timing. The cost to households who could not move is real. The institutional cash-flow cost is small only under the production floor form and face-value accounting, and at the zero-refinance edge of the rule-only band, an unrefereed import.
\end{abstract}

\vspace{0.5\baselineskip}
{\footnotesize\noindent\textbf{Keywords:} mortgage prepayment; lock-in effect; quantitative tightening; mortgage-backed securities; hazard models; monetary policy implementation.\\
\textbf{JEL codes:} E52, E58, G21, R31.\par}
\newpage